%-----------------------------------
\section{Estudo de caso}
\label{sec_tipo_estudo_caso}

\subsection{Caracterização}

O estudo de caso investiga um fenômeno contemporâneo em profundidade e em seu contexto real, especialmente quando os limites entre o fenômeno e o contexto não são claros \cite{yin2015estudo}. Em Sistemas de Informação, é comum estudar a implantação de um sistema em uma organização, diagnosticar e propor melhorias para processos de negócio \cite{dumas2018fundamentals} ou analisar como uma equipe adota uma prática de desenvolvimento \cite{runeson2009guidelines}.

O estudo de caso usa várias fontes de evidência, como entrevistas, documentos, observação e dados de sistemas, e as combina (triangulação) para fortalecer as conclusões. A autorização da organização estudada e o sigilo das informações sensíveis devem ser tratados na metodologia.

\subsection{Estrutura típica do trabalho}

\begin{quadro}[htb]
	\centering
	\caption{Estrutura típica de um TCC de estudo de caso.}
	\label{qua_estrutura_estudo_caso}
	\small
	\begin{tabularx}{\textwidth}{|l|X|}
		\hline
		\textbf{Capítulo} & \textbf{Seções típicas} \\ \hline
		Metodologia & Caracterização do estudo (caso único ou múltiplo, unidade de análise); fontes e técnicas de coleta; análise dos dados. \\ \hline
		O caso & Caracterização da organização; situação atual (AS-IS); proposta ou intervenção (TO-BE). \\ \hline
		Resultados e discussão & Resultados da análise ou da intervenção; discussão; limitações. \\ \hline
		Apêndices e anexos & Roteiro de entrevista (apêndice); documentos da organização (anexo). \\ \hline
	\end{tabularx}
	\fonte{Elaborado pelo autor.}
\end{quadro}

O Quadro~\ref{qua_estrutura_estudo_caso} resume a estrutura desse tipo de trabalho. Um modelo de roteiro de entrevista está no Apêndice~\ref{apD_roteiro}.

\subsection{Exemplo: processo atual e processo proposto}

Em estudos que envolvem processos de negócio, é comum comparar o processo atual (AS-IS) com o proposto (TO-BE), como na Figura~\ref{fig_processos}. Subfiguras permitem apresentar as duas versões lado a lado e referenciá-las individualmente: Figura~\ref{subfig_asis} e Figura~\ref{subfig_tobe}.

\begin{figure}[htb]
	\centering
	\caption{Processo de solicitação de compras de uma organização (exemplo).}
	\label{fig_processos}
	\tikzset{
		tarefa/.style={draw, rounded corners, align=center, font=\scriptsize, minimum width=2.4cm, minimum height=0.8cm},
		evento/.style={draw, circle, minimum size=0.45cm},
		fluxo/.style={-{Stealth}},
	}
	\begin{subfigure}{0.48\textwidth}
		\centering
		\begin{tikzpicture}[node distance=0.45cm]
			\node[evento] (i) {};
			\node[tarefa, below=of i] (t1) {Preencher\\formulário em papel};
			\node[tarefa, below=of t1] (t2) {Entregar ao setor\\de compras};
			\node[tarefa, below=of t2] (t3) {Digitar pedido\\na planilha};
			\node[tarefa, below=of t3] (t4) {Aprovar pedido};
			\node[evento, very thick, below=of t4] (f) {};
			\draw[fluxo] (i) -- (t1); \draw[fluxo] (t1) -- (t2); \draw[fluxo] (t2) -- (t3);
			\draw[fluxo] (t3) -- (t4); \draw[fluxo] (t4) -- (f);
		\end{tikzpicture}
		\caption{Processo atual (AS-IS).}
		\label{subfig_asis}
	\end{subfigure}
	\hfill
	\begin{subfigure}{0.48\textwidth}
		\centering
		\begin{tikzpicture}[node distance=0.45cm]
			\node[evento] (i) {};
			\node[tarefa, below=of i] (t1) {Registrar pedido\\no sistema};
			\node[tarefa, below=of t1] (t4) {Aprovar pedido\\no sistema};
			\node[evento, very thick, below=of t4] (f) {};
			\draw[fluxo] (i) -- (t1); \draw[fluxo] (t1) -- (t4); \draw[fluxo] (t4) -- (f);
		\end{tikzpicture}
		\caption{Processo proposto (TO-BE).}
		\label{subfig_tobe}
	\end{subfigure}
	\fonte{Elaborado pelo autor.}
\end{figure}
